\documentclass[12pt]{article}
\usepackage[utf8]{inputenc}
\usepackage[T2A]{fontenc}
\usepackage[russian]{babel}
\usepackage[a4paper,margin=2cm]{geometry}
\usepackage{hyperref}
\usepackage{enumitem}

\begin{document}

\begin{center}
	{\LARGE \textbf{Техническое задание}}\\[1em]
	{\Large Проект: Сервис для автоматического определения маркировки шины по фото}\\[2em]
	Версия документа: 1.0 \\
	Дата: 9 Мая, 2025 г.
\end{center}

\tableofcontents
\newpage

\section{Введение}
Цель данного документа — формализовать требования, архитектуру, сроки и критерии приёмки проекта системы для определения марки, модели и размерности шины по фото. Сервис предназначен для автоматического распознавания текста и параметров шины по изображению сбоку.

\section{Цели и задачи проекта}
\begin{itemize}[leftmargin=*]
	\item \textbf{Главная цель.} Принимать на вход изображение шины и выдавать:
	      \begin{itemize}
		      \item марку шины (из заранее определённого списка категорий);
		      \item модель шины (из заранее определённого списка категорий);
		      \item размерность (диаметр в единицах измерения, определённых заказчиком).
	      \end{itemize}
	\item \textbf{Ключевые метрики успеха (KPI).}
	      \begin{itemize}
		      \item Точность OCR (доля верно распознанных символов) - не менее 90\% корректных символов.
		      \item Время обработки одного запроса - не более 10 секунд.
		      \item Доля отказов - не более 1\%.
	      \end{itemize}
	\item \textbf{Заинтересованные стороны.}
	      \begin{itemize}
		      \item Заказчик (владелец автосервиса).
		      \item Разработчики и операционная команда.
	      \end{itemize}
\end{itemize}

\section{Функциональные требования}

\subsection{Входные данные}
\begin{itemize}[leftmargin=*]
	\item Форматы изображений: JPEG, PNG.
	\item Разрешение:
	      \begin{itemize}
		      \item Минимум: $512 \times 512$ пикселей.
		      \item Максимум: до $4096 \times 4096$. Перед обработкой изображение масштабируется (бикубическая интерполяция) к необходимому диапазону разрешений.
	      \end{itemize}
	\item Максимальный размер файла: по согласованию, не более 35 МБ.
\end{itemize}

\subsection{Выходные данные}
\begin{itemize}[leftmargin=*]
	\item \textbf{Марка} и \textbf{модель} шины — значения из заранее определённого каталога категорий.
	\item \textbf{Размерность} шины в формате, определяемым заказчиком.
	\item Формат ответа: текстовое сообщение с необходимыми заказчику данными.
\end{itemize}

\subsection{Критерии обработки неуверенных результатов}
\begin{itemize}[leftmargin=*]
	\item Если точность распознавания (заранее обсуждается с заказчиком) ниже порога, результат:
	      \begin{itemize}
		      \item отправляется на ручную проверку сотрудником, либо
		      \item клиенту возвращается ошибка с рекомендацией переснять фотографию.
	      \end{itemize}
\end{itemize}

\subsection{Пользовательские сценарии}
\begin{enumerate}[leftmargin=*]
	\item Сотрудник получает фотографию от пользователя.
	\item Сотрудник отправляет фото в Telegram-бот, который через REST API вызывает сервис.
	\item Сервис возвращает JSON-ответ с параметрами шины.
	\item Сотрудник проверяет результат и подтверждает его, либо валидирует.
\end{enumerate}

\section{Нефункциональные требования}

\subsection{Производительность}
\begin{itemize}[leftmargin=*]
	\item Время отклика: не более 10 секунд на обработку запроса.
	\item Пропускная способность: до 10 запросов в минуту (пик).
\end{itemize}

\subsection{Надёжность и доступность}
\begin{itemize}[leftmargin=*]
	\item Доступность сервиса (SLA): не менее 99\% uptime.
	\item Гарантийная поддержка: 6 месяцев с момента ввода в эксплуатацию.
\end{itemize}

\subsection{Масштабируемость}
\begin{itemize}[leftmargin=*]
	\item При необходимости масштабирования сервиса, продукт должен быть готов к вертикальному методу масштабируемости (увеличение вычислительной мощности на единственном сервере).
\end{itemize}

\subsection{Безопасность и приватность}
\begin{itemize}[leftmargin=*]
	\item TLS 1.3 для всех соединений.
	\item Данные на диске не сохраняются (кроме кода и временных артефактов обработки).
	\item Отсутствие хранения пользовательских метаданных.
\end{itemize}

\section{Архитектура и технологии}

\subsection{ML-компоненты}
\begin{itemize}[leftmargin=*]
	\item Текстовая детекция и OCR на базе новейших SOTA-решений.
	\item При необходимости модели дообучаются на специфичном датасете, предоставленном заказчиком.
\end{itemize}

\subsection{Данные и аннотация}
\begin{itemize}[leftmargin=*]
	\item Источники: публичные датасеты, размеченные фотографии, предоставленные заказчиком.
	\item Аннотация: ручная разметка сотрудником.
\end{itemize}

\subsection{Инфраструктура и контейнеризация}
\begin{itemize}[leftmargin=*]
	\item Сервис упакован в Docker-контейнер, запущен на защищённом хостинге.
	\item Доступ к сервису осуществляется через REST API с аутентификацией по API-ключу.
\end{itemize}

\subsection{Аппаратные требования}
\begin{itemize}[leftmargin=*]
	\item SSD: не менее 64 ГБ.
	\item Оперативная память: не менее 16 ГБ.
	\item CPU: не менее 4 ядер.
	\item GPU: видеопамять не менее 8 ГБ (рекомендуется серия NVIDIA RTX).
\end{itemize}

\section{Тестирование и валидация}
\begin{itemize}[leftmargin=*]
	\item Проверка корректности на валидационном наборе данных с известными метками при помощи тестов.
	\item Логирование каждого запроса: время, содержимое запроса, ответ сервиса.
\end{itemize}

\section{План внедрения и поддержки}

\subsection{Этапы разработки}
\begin{enumerate}[leftmargin=*]
	\item Предпроектное исследование и сбор датасета — до 2 недель.
	\item Разработка прототипа (тестовый запуск, без деплоя) — 2.5 недель.
	\item Запуск MVP (деплой, базовая функциональность) — 2.5 недель.
	\item Бета-тестирование (сотрудник и заказчик) — 2 недели.
	\item Исправление ошибок, подготовка продукта к деплою, вывод в продакшен — 3 недели.
\end{enumerate}
\textbf{Итого:} 12 недели (3 месяца), включая резерв (+1 месяц) на непредвиденные обстоятельства.


\subsection{Гарантийная поддержка}
\begin{itemize}[leftmargin=*]
	\item Сервис поддерживается разработчиком в течение 6 месяцев после ввода в эксплуатацию.
	\item В течение гарантийного срока - исправление ошибок, дополнительная валидация сервиса по запросу заказчика.
\end{itemize}

\section{Коммерческие условия}

\subsection{Бюджет проекта}
\begin{itemize}[leftmargin=*]
    \item Общая стоимость работ: 450000 руб.
    \item Стоимость включает в себя:
            \begin{itemize}
                \item Все этапы разработки, отладки, тестирования и поддержки в гарантийных случаях
                \item Все необходимые лицензии на ПО и инструменты, используемые в проекте, датасеты
                \item Все необходимые компоненты инфраструктуры на время разработки (сервер, хостинг, облачные сервисы)
                \item Оплату услуг специалиста, занимающегося исследованиями в области компьютерного зрения
            \end{itemize}
\end{itemize}

\subsection{Оплата}
\begin{itemize}[leftmargin=*]
    \item Оплата производится по факту сдачи проекта.
\end{itemize}

\subsection{Дополнительные расходы}
\begin{itemize}[leftmargin=*]
    \item При необходимости дополнительных задач или доработок, не включённых в базовое ТЗ, стоимость и сроки согласуются отдельно.
    \item Компоненты инфраструктуры (хостинг, лицензионное ПО, облачные сервисы) оплачивает заказчик по факту.
\end{itemize}

\section{Документация}
\begin{itemize}[leftmargin=*]
	\item Пользовательская документация с описанием API, примерами запросов/ответов и разбором типовых кейсов.
	\item Техническая документация для администратора (как запустить сервис).
\end{itemize}

\end{document}
